\section{Voting rules for a DIS}

A DIS is considered an enquiry draft and, consequently, enquiry stage rules apply.
According to the ISO/IEC Directives~\cite{jtc1:directives} (section 2.6.3),
the draft is approved if

\vspace{1em}
\begin{tcolorbox}[
  colback=blue!8!white,
  colframe=blue,
  title={\textbf{Section 2.6.3 from ISO/IEC JTC1 Directives}},
  breakable,
  width=\textwidth,
  boxrule=0.5pt
  ]
An enquiry draft is approved if

\begin{enumerate}[label=(\alph*)]

\item
a two-thirds majority of the votes cast by the P-members are in favour, and

\item
not more than one-quarter of the total number of votes cast are negative.

\end{enumerate}

Abstentions are excluded when the votes are counted, as well as negative votes
not accompanied by any comments of a technical nature. Comments received after
the normal voting period are submitted to the committee secretariat for
consideration at the time of the next review of the International Standard.
\end{tcolorbox}

\vspace{1em}
\begin{tcolorbox}[
  colback=red!8!white,
  colframe=red,
  title={\textbf{Results for DIS 14882 (Ed. 8)}},
  breakable,
  width=\textwidth,
  boxrule=0.5pt
  ]
In this specific case:

\begin{itemize}

\item
Criterion (a) received 12 of 20 P-member votes in favor (60\%), below the
two-thirds requirement (approximately 66.7\%).

\item
There were 9 negative votes among 22 Yes/No votes (approximately 40.9\%), above
the maximum of 25\%.
\end{itemize}

These calculations assume that all eight P-member No votes and all nine No
votes overall were accompanied by technical comments. Under section 2.6.3,
negative votes without such comments are excluded from the relevant count.
Either criterion being unmet is sufficient for disapproval.
\end{tcolorbox}

As a result of this outcome, section 2.6.4 needs to be applied:

\vspace{1em}
\begin{tcolorbox}[
  colback=blue!8!white,
  colframe=blue,
  title={\textbf{Section 2.6.4 from ISO/IEC JTC1 Directives}},
  breakable,
  width=\textwidth,
  boxrule=0.5pt
  ]
On receipt of the results of the voting and any comments, an appropriate group
of the committee, which would typically be the working group to which the
project is assigned, will, at a meeting, address any comments received and
prepare, as needed, a revised version. The group will advise the Chair on the
possible courses of action. [\ldots]
\end{tcolorbox}

Our interpretation is that WG21 should address all comments and prepare a
revised version of the draft. It has been suggested by Timur Doumler
and others~\cite{p4433r0} that WG21 should discuss and then dismiss all requests
for contracts removal. Literally it states the following: \emph{``Any comments
asking for the removal of contract assertions from the DIS should be discussed
(as procedurally required), but, absent new technical information, WG21 should
reject them as having no consensus for change and attach this paper to its
response as supporting reasoning.''}

However, the impact of such a decision should be carefully considered. Of the
34 recorded responses, 13 were Yes, 9 were No, and 12 were abstentions.
Abstentions are not votes against the DIS, and the tally alone does not show
whether contracts drove each response.
